% ===========================================================================
% Section: Multi-Modal Sensor Fusion
% ===========================================================================
\section{Multi-Modal Sensor Fusion}
\label{sec:sensor_fusion}

This section describes the two key components of the sensor fusion layer:
the pluggable \texttt{ISensorSource} interface and the multi-modal depth
fusion module.

\subsection{Pluggable Sensor Source Interface}
\label{subsec:sensor_interface}

The \texttt{ISensorSource} abstract interface (Figure~\ref{fig:sensor_architecture})
is the cornerstone of the framework's extensibility.

% ===========================================================================
% Figure: Sensor architecture — pluggable SDK interface
% Shows ISensorSource as an abstract interface with two concrete implementations.
% ===========================================================================
\begin{figure*}[t]
\centering
\begin{tikzpicture}[node distance=1.2cm]

% ---- Abstract interface (top) ----
\node[block, minimum width=6cm, minimum height=1.5cm, fill=blue!15]
  (iface) {\textbf{\texttt{ISensorSource}}\\[2pt]
           \texttt{+ open()} \quad \texttt{+ hasNext()}\\
           \texttt{+ nextFrame()} \quad \texttt{+ close()}};

% ---- Concrete implementations (bottom) ----
\node[blockgreen, below left=1.5cm and 0.5cm of iface, minimum width=4.5cm,
      minimum height=2.2cm]
  (reader) {\textbf{DatasetReader}\\[3pt]
            \footnotesize Reads from disk:\\
            \footnotesize • RGB images (JPG/PNG)\\
            \footnotesize • Depth maps (PNG/EXR)\\
            \footnotesize • LiDAR PLY files\\
            \footnotesize • Poses (TUM/KITTI/EuRoC)\\
            \footnotesize • IMU (CSV)};

\node[blockorange, below right=1.5cm and 0.5cm of iface, minimum width=4.5cm,
      minimum height=2.2cm]
  (orbbec) {\textbf{OrbbecSource}\\[3pt]
            \footnotesize Live camera stream:\\
            \footnotesize • OrbbecSDK v2 (C API)\\
            \footnotesize • Femto Mega / Gemini\\
            \footnotesize • LiDAR + iToF + stereo\\
            \footnotesize • IMU (on-board)\\
            \footnotesize \emph{(compile-gated)}};

% ---- Future extension ----
\node[block, below=3.5cm of iface, minimum width=4.5cm, fill=gray!10, dashed]
  (future) {\textbf{FutureSource} (extensible)\\[2pt]
            \footnotesize Any SDK implementing\\
            \footnotesize \texttt{ISensorSource}};

% ---- Inheritance arrows ----
\draw[dasharrow] (reader.north) -- (iface.south -| reader.north)
  node[midway, right, font=\footnotesize] {\texttt{implements}};
\draw[dasharrow] (orbbec.north) -- (iface.south -| orbbec.north)
  node[midway, left, font=\footnotesize] {\texttt{implements}};
\draw[dasharrow] (future.north) -- (iface.south)
  node[midway, right, font=\footnotesize] {\texttt{implements}};

% ---- Group ----
\begin{scope}[on background layer]
  \node[group=Sensor Source Layer, fit=(iface)(reader)(orbbec)(future),
        inner sep=14pt, fill=blue!2] {};
\end{scope}

\end{tikzpicture}
\caption{The \texttt{ISensorSource} abstract interface and its concrete
implementations. The \texttt{DatasetReader} handles offline file-based data,
while the \texttt{OrbbecSource} wraps the Orbbec SDK for live camera input.
The interface is designed for extensibility: any future sensor SDK can be
integrated by implementing four virtual methods.}
\label{fig:sensor_architecture}
\end{figure*}


It defines four virtual methods:

\begin{itemize}
    \item \texttt{open(config)}: Initialise the source from a YAML
    configuration node. For offline readers, this opens files and scans
    directories; for live cameras, this connects to the device and starts
    streaming.

    \item \texttt{hasNext()}: Poll whether more frames are available.
    For offline readers, this checks the frame index against the total
    count; for live cameras, this returns \texttt{true} until the stream
    is stopped.

    \item \texttt{nextFrame()}: Retrieve the next \texttt{AlignedFrame}.
    For offline readers, this loads the next set of files from disk;
    for live cameras, this blocks until the next frame set arrives.

    \item \texttt{close()}: Release resources.
\end{itemize}

\subsubsection{DatasetReader}

The \texttt{DatasetReader} is the primary sensor source for offline
processing. It reads from a structured directory layout:

\begin{verbatim}
base_path/
  images/         RGB frames (000000.jpg, ...)
  depth/          Depth maps (000000.png [16-bit mm] or .exr)
  pointclouds/    LiDAR PLY files (000000.ply, ...)
  trajectory.txt  Camera poses (TUM/KITTI/EuRoC format)
  imu.csv         IMU data (optional)
\end{verbatim}

Three pose formats are supported:
\begin{itemize}
    \item \textbf{TUM}: \texttt{timestamp tx ty tz qx qy qz qw}
    \item \textbf{KITTI}: 3$\times$4 row-major rigid transform matrix
    ($R | \mathbf{t}$), no timestamp
    \item \textbf{EuRoC}: CSV with nanosecond timestamps
    (\texttt{timestamp\_ns, px, py, pz, qw, qx, qy, qz})
\end{itemize}

Depth maps are loaded as 32-bit float images in metric units. PNG files
store depth in millimetres (16-bit unsigned) and are automatically
converted; EXR files store metric depth directly.

Point clouds are read via a minimal PLY parser that supports both ASCII
and binary little-endian formats, handling \texttt{float32} XYZ coordinates
and \texttt{uchar} RGB colour.

\subsubsection{OrbbecSource}

The \texttt{OrbbecSource} wraps the OrbbecSDK v2 C API for live data
acquisition from Orbbec 3D cameras (Femto Mega, Gemini 2). It enables
three streams---colour (RGB), depth (iToF), and IR/LiDAR---and returns
\texttt{AlignedFrame} objects with the same structure as the offline
reader.

This implementation is \emph{compile-gated}: when the
\texttt{USE\_ORBBEC\_SDK} preprocessor macro is not defined, a stub
implementation is compiled that prints a helpful error message and
returns \texttt{false} from \texttt{open()}. This ensures the framework
compiles cleanly without the SDK, and users can enable live camera
support by installing the SDK and passing
\texttt{-DUSE\_ORBBEC\_SDK=ON} to CMake.

\subsection{Multi-Modal Depth Fusion}
\label{subsec:depth_fusion}

The depth fusion module (Figure~\ref{fig:depth_fusion}) combines three
depth sources into a single high-quality depth map.

% ===========================================================================
% Figure: Multi-modal depth fusion data flow
% LiDAR sparse + iToF dense + stereo SGM → confidence weighting → SPNet
% ===========================================================================
\begin{figure*}[t]
\centering
\begin{tikzpicture}[node distance=0.7cm and 1.2cm]

% ---- Row 1: Three depth sources ----
\node[blockred, minimum width=3cm] (lidar_d) {LiDAR Depth\\(sparse, $D_L$)};
\node[blockorange, right=of lidar_d, minimum width=3cm] (itof_d)
  {iToF Depth\\(dense, $D_F$)};
\node[blockgreen, right=of itof_d, minimum width=3cm] (stereo_d)
  {Stereo Depth\\(SGM, $D_S$)};

% ---- Row 2: Confidence maps ----
\node[block, below=of lidar_d, minimum width=3cm, fill=red!5]
  (conf_l) {$C_L$: binary\\(1.0 at hits)};
\node[block, below=of itof_d, minimum width=3cm, fill=orange!5]
  (conf_f) {$C_F$: edge-based\\(1.0 $- \nabla D$/max)};
\node[block, below=of stereo_d, minimum width=3cm, fill=green!5]
  (conf_s) {$C_S$: texture\\(variance-based)};

% ---- Row 3: Fusion node ----
\node[block, below=2.0cm of conf_f, minimum width=5cm, minimum height=1.2cm,
      fill=purple!12]
  (fuse) {Confidence-Weighted Fusion\\
          $D_{\text{fused}} = \frac{\sum_i C_i \cdot D_i}{\sum_i C_i}$};

% ---- Row 4: LiDAR override ----
\node[blockred, below=0.7cm of fuse, minimum width=5cm]
  (override) {LiDAR Override\\(ground-truth pixels)};

% ---- Row 5: SPNet completion ----
\node[blockorange, below=0.7cm of override, minimum width=5cm]
  (spnet) {SPNet TensorRT\\(hole completion)};

% ---- Row 6: Final depth ----
\node[blockgreen, below=0.7cm of spnet, minimum width=5cm]
  (final) {Fused + Completed Depth\\$D_{\text{final}}$};

% ---- Arrows: sources → confidence ----
\draw[arrow] (lidar_d) -- (conf_l);
\draw[arrow] (itof_d)  -- (conf_f);
\draw[arrow] (stereo_d)-- (conf_s);

% ---- Arrows: confidence → fusion ----
\draw[arrow] (conf_l.south) -- (fuse.north -| conf_l.south);
\draw[arrow] (conf_f.south) -- (fuse.north);
\draw[arrow] (conf_s.south) -- (fuse.north -| conf_s.south);

% ---- Main flow ----
\draw[arrow] (fuse) -- (override);
\draw[arrow] (override) -- (spnet);
\draw[arrow] (spnet) -- (final);

% ---- Annotations ----
\node[font=\footnotesize, right=0.3cm of fuse, text=purple!70!black]
  {weighted average};
\node[font=\footnotesize, right=0.3cm of override, text=red!70!black]
  {$D_L$ overwrites where available};

\end{tikzpicture}
\caption{Multi-modal depth fusion pipeline. Three depth sources---sparse
LiDAR, dense iToF, and stereo SGM---are each assigned a confidence map.
A weighted average produces the fused depth, after which LiDAR pixels
overwrite the fused result as ground truth. The SPNet TensorRT engine
completes remaining holes.}
\label{fig:depth_fusion}
\end{figure*}


Each source has
complementary characteristics:

\begin{table}[h]
\centering
\caption{Depth source characteristics.}
\label{tab:depth_sources}
\begin{tabular}{lccc}
\toprule
Source & Density & Accuracy & Limitation \\
\midrule
LiDAR  & Sparse ($\sim$5\% pixels) & High ($\pm$2\,cm) & Limited spatial coverage \\
iToF   & Dense (100\%) & Medium ($\pm$5\,cm) & Noise at discontinuities \\
Stereo & Dense (100\%) & Medium ($\pm$10\,cm) & Fails in textureless areas \\
\bottomrule
\end{tabular}
\end{table}

\subsubsection{LiDAR Depth Projection}

LiDAR points in the world frame are projected into the camera image to
produce a sparse depth map. Given a point $\mathbf{p}_w$ and camera pose
$(R_{wc}, \mathbf{t}_{wc})$:

\begin{align}
    \mathbf{p}_c &= R_{cw} \, \mathbf{p}_w + \mathbf{t}_{cw} \\
    u &= f_x \, p_{c,x} / p_{c,z} + c_x \\
    v &= f_y \, p_{c,y} / p_{c,z} + c_y
\end{align}

A Z-buffer ensures that the nearest depth is retained when multiple
points project to the same pixel.

\subsubsection{Stereo Depth via SGM}

When stereo image pairs are available, depth is computed via OpenCV's
\texttt{StereoSGBM} (Semi-Global Block Matching). The disparity map $d(u,v)$
is converted to metric depth:

\begin{equation}
    Z(u,v) = \frac{b \cdot f_x}{d(u,v)}
\end{equation}

where $b$ is the stereo baseline and $f_x$ is the horizontal focal length.

\subsubsection{Confidence-Weighted Fusion}

Each depth source is assigned a per-pixel confidence weight $C_i(u,v)$:

\begin{itemize}
    \item \textbf{LiDAR} ($C_L$): Binary---1.0 at pixels with LiDAR hits,
    0 elsewhere. LiDAR depth is treated as ground truth.

    \item \textbf{iToF} ($C_F$): Edge-based confidence, reduced at depth
    discontinuities:
    $C_F = 1 - \|\nabla D_F\| / \max\|\nabla D_F\|$.

    \item \textbf{Stereo} ($C_S$): Texture-based confidence, proportional
    to local depth variance (a proxy for scene texture).
\end{itemize}

The fused depth is computed as a weighted average:
\begin{equation}
    D_{\text{fused}}(u,v) = \frac{\sum_{i \in \{L, F, S\}} C_i(u,v) \cdot D_i(u,v)}
                                  {\sum_{i \in \{L, F, S\}} C_i(u,v)}
\end{equation}

After fusion, LiDAR pixels overwrite the fused depth as ground truth:
$D_{\text{fused}}(u,v) \leftarrow D_L(u,v)$ wherever $D_L(u,v) > 0$.

\subsection{Timestamp-Based Data Alignment}

Figure~\ref{fig:data_alignment} illustrates the timestamp-based alignment
strategy.

% ===========================================================================
% Figure: Timestamp-based multi-modal data alignment
% Timeline showing synchronization across 4 modalities.
% ===========================================================================
\begin{figure*}[t]
\centering
\begin{tikzpicture}[node distance=0.5cm, font=\small]

% ---- Timeline axis ----
\draw[->, thick] (0,0) -- (13,0) node[right] {time (s)};
\foreach \x/\label in {0/0.0, 2/0.1, 4/0.2, 6/0.3, 8/0.4, 10/0.5, 12/0.6}
{
  \draw[thick] (\x, -0.1) -- (\x, 0.1);
  \node[below, font=\footnotesize] at (\x, -0.1) {\label};
}

% ---- LiDAR row ----
\node[left, font=\footnotesize\bfseries] at (-0.3, 1.5) {LiDAR};
\foreach \x in {0.5, 2.5, 4.5, 6.5, 8.5, 10.5}
{
  \draw[arrow, red!70!black] (\x, 0.8) -- (\x, 1.5);
  \node[font=\tiny, above] at (\x, 1.5) {$\PointCloud$};
}
\draw[red!70!black, thick, dashed] (0, 1.5) -- (12.5, 1.5);

% ---- iToF row ----
\node[left, font=\footnotesize\bfseries] at (-0.3, 2.8) {iToF};
\foreach \x in {0.3, 1.3, 2.3, 3.3, 4.3, 5.3, 6.3, 7.3, 8.3, 9.3, 10.3, 11.3}
{
  \draw[arrow, orange!80!black] (\x, 2.1) -- (\x, 2.8);
}
\draw[orange!80!black, thick, dashed] (0, 2.8) -- (12.5, 2.8);
\node[font=\tiny, right] at (12.5, 2.8) {$D_F$ (30 Hz)};

% ---- Stereo/RGB row ----
\node[left, font=\footnotesize\bfseries] at (-0.3, 4.1) {Stereo/RGB};
\foreach \x in {0.4, 1.4, 2.4, 3.4, 4.4, 5.4, 6.4, 7.4, 8.4, 9.4, 10.4, 11.4}
{
  \draw[arrow, green!60!black] (\x, 3.4) -- (\x, 4.1);
}
\draw[green!60!black, thick, dashed] (0, 4.1) -- (12.5, 4.1);
\node[font=\tiny, right] at (12.5, 4.1) {$\Image$ (30 Hz)};

% ---- IMU row ----
\node[left, font=\footnotesize\bfseries] at (-0.3, 5.4) {IMU};
\foreach \x in {0.05, 0.25, 0.45, 0.65, 0.85, 1.05, 1.25, 1.45, 1.65, 1.85,
                2.05, 2.25, 2.45, 2.65, 2.85, 3.05, 3.25, 3.45, 3.65, 3.85,
                4.05, 4.25, 4.45, 4.65, 4.85, 5.05, 5.25, 5.45, 5.65, 5.85,
                6.05, 6.25, 6.45, 6.65, 6.85, 7.05, 7.25, 7.45, 7.65, 7.85,
                8.05, 8.25, 8.45, 8.65, 8.85, 9.05, 9.25, 9.45, 9.65, 9.85,
                10.05, 10.25, 10.45, 10.65, 10.85, 11.05, 11.25, 11.45, 11.65, 11.85}
{
  \draw[blue!70!black] (\x, 4.9) -- (\x, 5.4);
}
\draw[blue!70!black, thick, dashed] (0, 5.4) -- (12.5, 5.4);
\node[font=\tiny, right] at (12.5, 5.4) {IMU (200 Hz)};

% ---- Aligned frame markers ----
\foreach \x/\label in {2.5/$t_1$, 6.5/$t_2$, 10.5/$t_3$}
{
  \draw[thick, purple!70!black, dashed] (\x, 0) -- (\x, 5.8);
  \node[font=\footnotesize, above, purple!70!black] at (\x, 5.8) {Aligned \label};
}

% ---- Tolerance band ----
\draw[<->, thick, gray] (2.4, 0.5) -- (2.6, 0.5)
  node[midway, below, font=\tiny] {$\pm\epsilon$};

\end{tikzpicture}
\caption{Timestamp-based multi-modal data alignment. The LiDAR stream
(10--20\,Hz) defines the reference timestamps. iToF, stereo, and IMU data
are matched within a tolerance $\epsilon = 10$\,ms. Dashed vertical lines
mark aligned frame instances $t_1, t_2, t_3$ where all modalities are
synchronized.}
\label{fig:data_alignment}
\end{figure*}


The LiDAR stream (10--20\,Hz) defines the reference timestamps,
and other modalities are matched within a tolerance
$\epsilon = 10$\,ms. This mirrors the alignment logic in the original
\gaussianlic{} system, where ROS message timestamps were used for
queue-based synchronization.

In the offline \texttt{DatasetReader}, alignment is implicit because all
modalities share the same frame numbering. The \texttt{DataAligner}
class provides an explicit alignment utility for live streaming scenarios
where modalities arrive at different rates.
